% Pursuing Stacks (cotes 134-2 to 134-8): archive cote 134-7, site page 60 (PDF page 61), typescript page 502.
% typescript. Draft compared: ps6.tex:2948-3007.
% AI-transcribed by Claude Opus 5.5 (claude-opus-5-5, adaptive thinking, effort high), Message Batches API, batch ps4 (msgbatch_012Aoc1NuskATPsyBeKSYK7J), 2026-11-12; not yet checked by eye.
\scanpage{134-7/60}
\mspage{502}

or successive extensions of such) to give the same result.  Namely \deleted{the}
the latter is computed in terms of cochains which are polynomial functi\over{on}{on}s
\uline{with coefficients in $\underline{Z}$} ,  whereas there exist polynomial functtions \uline{with}
\uline{coefficients in $\underline{Q}$} \inserted{(not in $\underline{Z}$)} which\hand{,} however\hand{,} give rise to integer-valued functions on
the group of integer-valued points\over{,}{.} \deleted{and notably} \inserted{Such are} the binomial expressions
\[ P_n(t) = t(t-1)\dots(t-n+1)/n! \qquad \hand{(for $n \in \underline{N}$).} \]

These (in the case of just one variable t) are known to form a basis of
the $\underline{Z}$-module of all integer-valued functions on $\underline{Z}$ ,  and there is a corres-%
ponding basis for integer-valued functions on $\underline{Z}^r$\hand{,} for any natural integer
r.  Thus the hope still remains that a sweeping comparison theorem for
discrete versus ``schematic'' \xout{(afxsertx)} linearization might hold true,
provided it is expressed in such a way that the ``schematic models'' we are
working with should be built up with ``schemes'' (of sorts) described in
terms of spectra not of polynomial algebra\unsure{s}  $\underline{Z}\hand{[}t\hand{]}$ and tensor powers of
these,  but rather of  ``binomial algebra\unsure{s}''  $\underline{Z}\hand{$\langle t\rangle$}$ built up with the binomial
expressions above,  and tensor powers  of such.  If we want to develop a
corresponding notion of homotopy types over a general gound ring k,  we
should then require upon k an extra structure of a ``binomial ring'' (as
introduced in the Riemann-Roch Seminar SGA 6 in some talks of Berthelot),
namely a ring endowed with operations
\[ x \hand{$\mapsto$} \binom{x}{n} \ : \ k \hand{$\to$} k \qquad (n \hand{$\in$} \hand{$\underline{N}$}) \ , \]
satisfying the formal properties of the binomial functions  $x \hand{$\mapsto$} P_n(x)$
in the case $k=\underline{Z}$ or $\underline{Q}$ .  Whereas linearization of homotopy types via De
Rham complexes with devided powers relies on a ``commutative algebra with
d\over{\unsure{e}}{e}vided powers'' \xout{\unsure{and corresponding}} \hand{(which} was developped extensively by Berthe-%
lot and others for the needs of crystalline cohomology\hand{)}\hand{,}  linearization via
unipotent bundles (assuming it can be done in such a way as to insure that
any discrete homotopy type can be``schematized''in an essentially unique way\hand{)}
might well rely on the development of a ``binomial commutative algebra''
and a corresponding notion of ``binomial schemes''.  There should be a lot
of fun ahead developping the necessary algebraic machinery,  which may
prove of interest in it's own right.  It should be realized however that \margin{See comments next section (p.506-507)}
for a ring k to admit a binomial structure is a rather strong restriction
\xout{sxxk}  -  thus for a given prime p,  no field of char.p (except possibly
the prime field ?) admits such a structure.  This remark may temper some-%
what the enthousiasm for pushing in this direction, \inserted{even} granting that a
``binomial comparison theorem'' for di\deleted{\unsure{t|e}}screte versus ``binomial'' linearizati-%
on holds true.

\margin{60}
